% Class-compatibility preamble retained for a reader excerpt of
%   Lisa Sauermann, Dmitrii Zakharov, "A Sharp Ramsey Theorem for Ordered
%   Hypergraph Matchings", arXiv:2309.04813v4, CC BY 4.0.
%
% The source pins the publisher class aic (the AIMS journal class), which is not part of
% the delivery baseline, so this file supplies the amsart baseline and then copies the
% author's own declaration block mechanically and unchanged
% (pinned-original-source.tex lines 25 and 28-59: the ams packages, the author's eleven
% newtheorem declarations with [theorem] as the shared counter, the definition theorem
% style, the author's comment that the file should add as few macro or package
% definitions as possible, and the author's own macros \su, \sm, \Z, \HH, \conv and \PP).
% The author's tikz package (pinned source line 26) serves the three figures of the paper
% that this excerpt does not carry and is therefore not reproduced; the author's babel,
% fontenc and inputenc lines (pinned source lines 47-49) are layout/font choices of the
% publisher build and are not part of the mathematical excerpt, the delivery engine being
% XeLaTeX with UTF-8 input.
% No mathematical macro, statement, symbol or assumption is introduced, changed or
% removed by this file.
\documentclass[10pt,a4paper]{amsart}
\usepackage{amsthm,amsfonts,amssymb,amsmath}
\usepackage{graphicx}

\newtheorem{theorem}{Theorem}[section]
\newtheorem{proposition}[theorem]{Proposition}
\newtheorem{lemma}[theorem]{Lemma}
\newtheorem{conjecture}[theorem]{Conjecture}
\newtheorem{fact}[theorem]{Fact}
\newtheorem{claim}[theorem]{Claim}
\newtheorem{corollary}[theorem]{Corollary}
\newtheorem{definition}[theorem]{Definition}
\newtheorem{problem}[theorem]{Problem}
\newtheorem{question}[theorem]{Question}
\newtheorem{observation}[theorem]{Observation}

\theoremstyle{definition}
\newtheorem{remark}[theorem]{Remark}

% The author notes at this point that as few macro or package
% definitions as possible should be added; the six macros below are the
% author's own and are copied unchanged.

\newcommand{\su}{\subseteq}
\newcommand{\sm}{\setminus}
\newcommand{\Z}{\mathbb{Z}}
\newcommand{\HH}{\mathcal{H}}
\newcommand{\conv}{\operatorname{conv}}
\newcommand{\PP}{\mathcal{P}}
